\section*{Chapter \ref{ch:b1:layered-earth} --- The Layered Earth}

\begin{solution}{exo:b1:layered-earth:1}
The volume fractions are the cubes of the radius ratios: core
$(3480/6371)^3 = 0.1630$, that is \qty{16.3}{\percent}; \omterm{def:b1:layered-earth:layers}{inner core}
$(1221.5/6371)^3 = 0.00705$, \qty{0.70}{\percent}.
\end{solution}

\begin{solution}{exo:b1:layered-earth:2}
Above: $\Phi = 13.7166^2 - \tfrac43\times7.26466^2 = 188.145 - 70.367 =
\qty{117.78}{km^2.s^{-2}}$, $K = \rho\Phi = \qty{5566.45}{kg.m^{-3}}\times
\qty{117.78e6}{m^2.s^{-2}} = \qty{655.6}{GPa}$. Below: $\Phi = 8.06482^2 =
\qty{65.04}{km^2.s^{-2}}$, $K = 9903.49\times\num{65.04e6} = \qty{644.1}{GPa}$.
The bulk modulus hardly changes (\qty{-1.7}{\percent}). What changes is the
shear modulus, $\mu = \rho v_S^2 = \qty{293.8}{GPa}$ above and $0$ below (a
liquid), and the density, which rises by \qty{78}{\percent}: hence the drop of
$v_P$ despite an almost equal $K$.
\end{solution}

\begin{solution}{exo:b1:layered-earth:3}
$\sin i_c = 6.2/8.1 = 0.7654$, $i_c = \ang{49.95}$. Critical distance
$2\times30\times\tan\ang{49.95} = 60\times1.1906 = \qty{71.4}{km}$. \omterm{def:b1:layered-earth:head-wave}{Crossover
distance} $60\sqrt{14.3/1.9} = 60\times2.743 = \qty{164.6}{km}$.
\end{solution}

\begin{solution}{exo:b1:layered-earth:4}
$g = Gm/r^2 = \num{6.674e-11}\times\num{1.94e24}/(\num{3.48e6})^2 =
\qty{10.69}{m.s^{-2}}$, \qty{9}{\percent} more than at the surface. Going down
to the core, the radius squared falls to $(3480/6371)^2 = 0.298$ of its surface
value, while the attracting mass falls only to $1.94/5.972 = 0.325$ of the total,
because the core holds a third of the mass: $0.325/0.298 = 1.09$.
\end{solution}

\begin{solution}{exo:b1:layered-earth:5}
Near the source $v_1 = 60/10.0 = \qty{6.0}{km.s^{-1}}$. Far away
$v_2 = (300 - 200)/(45.8 - 33.3) = 100/12.5 = \qty{8.0}{km.s^{-1}}$; the
intercept time is $t_i = 33.3 - 200/8.0 = \qty{8.3}{s}$. Then $\sin i_c =
0.75$, $\cos i_c = 0.6614$, and $h = t_i v_1/(2\cos i_c) = 8.3\times6.0/1.3229
= \qty{37.6}{km}$. Check: $x_c = 2\times37.6\times\sqrt{14/2} = \qty{199}{km}$,
where both waves arrive at $199/6.0 = 199/8.0 + 8.3 = \qty{33.2}{s}$.
\end{solution}

\begin{solution}{exo:b1:layered-earth:6}
For uniform $\rho$, $I = \tfrac{8\pi}{3}\rho\int_0^R r^4\,\mathrm dr =
\tfrac{8\pi}{15}\rho R^5$ and $M = \tfrac43\pi\rho R^3$, so
$I/MR^2 = (8/15)/(4/3) = 0.4$. Two layers ($r_c = \qty{3.48e6}{m}$,
$R = \qty{6.371e6}{m}$): $M = \tfrac43\pi\bigl(11000\,r_c^3 + 4450\,(R^3 -
r_c^3)\bigr) = \qty{5.977e24}{kg}$, \qty{0.08}{\percent} above the measured
mass; $I = \tfrac{8\pi}{15}\bigl(11000\,r_c^5 + 4450\,(R^5 - r_c^5)\bigr)$ gives
$I/MR^2 = 0.3457$. The mass is right, but the factor is \qty{4.5}{\percent}
too high: the real \omterm{def:b1:layered-earth:layers}{mantle} is not uniform, its density grows with depth, which
moves mass inward.
\end{solution}

\begin{solution}{exo:b1:layered-earth:7}
$\Phi = 12.16^2 - \tfrac43\times6.673^2 = 147.87 - 59.37 =
\qty{88.49}{km^2.s^{-2}}$. With depth $z$,
$\mathrm d\rho/\mathrm dz = \rho g/\Phi = 4860\times9.9/\num{88.49e6} =
\qty{5.44e-4}{kg.m^{-4}}$, that is \qty{0.544}{kg.m^{-3}} per kilometre, or
\qty{0.54}{g.cm^{-3}} per \qty{1000}{km}. PREM rises by $5.357 - 4.844 =
0.513$ over these \qty{1000}{km}, \qty{6}{\percent} less: the gradient
computed at the top is applied to the whole interval, while $\Phi$ grows with
depth and the gradient decreases.
\end{solution}

\begin{solution}{exo:b1:layered-earth:8}
$P \approx \bar\rho g z = 3570\times9.9\times\num{670e3} = \qty{23.7}{GPa}$
(PREM: \qty{23.9}{GPa}). It is the pressure, about \qty{24}{GPa}, at which
ringwoodite breaks down into bridgmanite and ferropericlase: the
\qty{660}{km} discontinuity.
\end{solution}

\begin{solution}{exo:b1:layered-earth:9}
(a)~Liquid iron: $10.90\times1.05 = 11.4$ to $10.90\times1.10 =
\qty{12.0}{g.cm^{-3}}$. (b)~$12.764/12.166 = 1.049$: a jump of
\qty{4.9}{\percent} in PREM, a little more than the \qty{4.5}{\percent} of
recent studies. (c)~Freezing alone would make the solid only slightly denser
than the liquid; the larger jump means that the solid holds fewer light
elements than the liquid: they prefer the liquid, and freezing at the
\omterm{def:b1:layered-earth:discontinuity}{inner-core boundary} leaves most of them behind.
\end{solution}

\begin{solution}{exo:b1:layered-earth:10}
Let $x_* = 2h\tan i_c$. Then $\sqrt{x_*^2 + 4h^2} = 2h/\cos i_c$ and
$t_r(x_*) = 2h/(v_1\cos i_c)$. The \omterm{def:b1:layered-earth:head-wave}{head wave} gives
$t(x_*) = 2h\tan i_c/v_2 + 2h\cos i_c/v_1 = \frac{2h}{v_1}\bigl(\sin^2 i_c/\cos
i_c + \cos i_c\bigr) = 2h/(v_1\cos i_c)$, using $\sin i_c = v_1/v_2$: the same.
The slope of the hyperbola is $t_r'(x) = x/(v_1\sqrt{x^2 + 4h^2})$; at $x_*$ it is
$\sin i_c/v_1 = 1/v_2$, the slope of the head-wave line. Since
$t_r'' > 0$, the hyperbola is convex and lies above each of its tangents: for
$x > x_*$, $t_r > t_{\text{head}}$. At the critical distance the reflection
strikes the boundary at exactly $i_c$ and is the \omterm{def:b1:layered-earth:head-wave}{head wave}'s birth; beyond,
the \omterm{def:b1:layered-earth:head-wave}{head wave}, which gains by running in the fast medium, comes first.
\end{solution}

\begin{solution}{exo:b1:layered-earth:11}
$m(r) = \tfrac43\pi\rho r^3$, so $g = Gm/r^2 = \tfrac43\pi G\rho r$. Then
$P(r) = \int_r^R \rho\,\tfrac43\pi G\rho s\,\mathrm ds = \tfrac23\pi G\rho^2(R^2
- r^2)$. At the centre $P = \tfrac23\pi\times\num{6.674e-11}\times5513^2\times
(\num{6.371e6})^2 = \qty{172}{GPa}$, \qty{47}{\percent} of PREM's
\qty{363.85}{GPa}: concentrating the mass towards the centre roughly doubles
the central pressure.
\end{solution}

\begin{solution}{exo:b1:layered-earth:12}
With $z$ the depth, $\mathrm d\rho/\mathrm dz = \rho g/\Phi$; if $g/\Phi$ is a
constant $a$, $\mathrm d(\ln\rho)/\mathrm dz = a$ and $\rho = \rho_0
e^{a(z - z_0)}$. Here $a(z - z_0) = 7.9\times\num{2258.5e3}/\num{87e6} =
0.2051$ and $\rho = 9.90\times e^{0.2051} = 9.90\times1.2276 =
\qty{12.15}{g.cm^{-3}}$, \qty{0.2}{\percent} below PREM's $12.17$: the whole
increase of density through the \omterm{def:b1:layered-earth:layers}{outer core} is self-compression.
\end{solution}

\begin{solution}{pb:b1:layered-earth:1}
\textbf{1.} $V = \tfrac43\pi(\num{6.371e6})^3 = \qty{1.0832e21}{m^3}$.

\textbf{2.} $\tfrac43\pi(r_2^3 - r_1^3)$: \omterm{def:b1:layered-earth:layers}{crust}, upper \omterm{def:b1:layered-earth:layers}{mantle} and transition zone
\qty{3.0706e20}{m^3}; lower \omterm{def:b1:layered-earth:layers}{mantle} \qty{5.9961e20}{m^3}; \omterm{def:b1:layered-earth:layers}{outer core}
\qty{1.6890e20}{m^3}; \omterm{def:b1:layered-earth:layers}{inner core} \qty{7.634e18}{m^3}.

\textbf{3.} \qty{28.35}{\percent}, \qty{55.36}{\percent}, \qty{15.59}{\percent}
and \qty{0.70}{\percent}: the lower \omterm{def:b1:layered-earth:layers}{mantle} is the largest shell, more than half
of the Earth's volume.

\textbf{4.} $15.59 + 0.70 = \qty{16.3}{\percent}$.

\textbf{5.} \omterm{def:b1:layered-earth:layers}{Crust}, upper \omterm{def:b1:layered-earth:layers}{mantle} and transition zone $\num{3.0706e20}\times3570 =
\qty{1.0962e24}{kg}$; lower \omterm{def:b1:layered-earth:layers}{mantle} \qty{2.9381e24}{kg}; \omterm{def:b1:layered-earth:layers}{outer core}
\qty{1.8410e24}{kg}; \omterm{def:b1:layered-earth:layers}{inner core} \qty{9.841e22}{kg}.

\textbf{6.} \qty{5.9737e24}{kg}, \qty{0.03}{\percent} above the measured
\qty{5.972e24}{kg}.

\textbf{7.} $\num{5.9737e24}/\num{1.0832e21} = \qty{5515}{kg.m^{-3}}$.

\textbf{8.} $\num{9.841e22}/\num{5.9737e24} = \qty{1.65}{\percent}$.

\textbf{9.} $(1.0962 + 2.9381)/5.9737 = \qty{67.5}{\percent}$.

\textbf{10.} The core is \qty{16.3}{\percent} of the volume and
\qty{32.5}{\percent} of the mass, twice as much: its mean density,
\qty{10986}{kg.m^{-3}}, is about $2.5$ times that of the \omterm{def:b1:layered-earth:layers}{mantle}.

\textbf{11.} A shell of uniform density gives
$I = \tfrac{8\pi}{3}\rho\int_{r_1}^{r_2} r^4\,\mathrm dr = \tfrac{8\pi}{15}\rho
(r_2^5 - r_1^5)$.

\textbf{12.} \omterm{def:b1:layered-earth:layers}{Crust}, upper \omterm{def:b1:layered-earth:layers}{mantle} and transition zone \qty{2.676e37}{kg.m^2}; lower \omterm{def:b1:layered-earth:layers}{mantle}
\qty{4.525e37}{kg.m^2}; \omterm{def:b1:layered-earth:layers}{outer core} \qty{9.27e36}{kg.m^2}; \omterm{def:b1:layered-earth:layers}{inner core}
\qty{5.87e34}{kg.m^2}.

\textbf{13.} The lower \omterm{def:b1:layered-earth:layers}{mantle}: it has half the mass, at large radii. The \omterm{def:b1:layered-earth:layers}{inner
core} holds \qty{1.6}{\percent} of the mass, and all of it close to the axis:
its share is \qty{0.07}{\percent}.

\textbf{14.} $I = \qty{8.1345e37}{kg.m^2}$ and $I/MR^2 = \num{8.1345e37}/
(\num{5.9737e24}\times(\num{6.371e6})^2) = 0.3355$.

\textbf{15.} \qty{1.4}{\percent} above $0.3308$, and well below $0.4$: the
model has the right concentration of mass towards the centre, nearly.

\textbf{16.} Yes. Inside each shell the real density is higher at the bottom
than at the top; a uniform shell of the same mass puts too much of that mass
near its outer radius, where it counts with $r^2$, and so overestimates $I$.

\textbf{17.} $g = GM/R^2 = \num{6.674e-11}\times\num{5.9737e24}/
(\num{6.371e6})^2 = \qty{9.822}{m.s^{-2}}$.

\textbf{18.} The core's mass is $\num{1.8410e24} + \num{9.841e22} =
\qty{1.9394e24}{kg}$, and $g = G m/r_c^2 = \qty{10.69}{m.s^{-2}}$.

\textbf{19.} $(\num{1.0962e24} + \num{2.9381e24})/(\num{3.0706e20} +
\num{5.9961e20}) = \num{4.0343e24}/\num{9.0667e20} = \qty{4450}{kg.m^{-3}}$.

\textbf{20.} $P \approx 4450\times10\times\num{2.891e6} = \qty{128.6}{GPa}$,
\qty{5}{\percent} below PREM's \qty{135.75}{GPa}.

\textbf{21.} $\rho_c = \num{1.9394e24}/\num{1.7653e20} = \qty{10986}{kg.m^{-3}}$,
and $\Delta P = \tfrac23\pi\times\num{6.674e-11}\times10986^2\times
(\num{3.48e6})^2 = \qty{204.3}{GPa}$.

\textbf{22.} $128.6 + 204.3 = \qty{332.9}{GPa}$, \qty{8.5}{\percent} below
PREM's \qty{363.85}{GPa}. Both estimates ignore the increase of density with
depth inside each region, which places more mass deeper, keeps the gravity
higher on the way down and raises $\rho g$, the pressure gradient.

\textbf{23.} The total mass (question 6), the \omterm{def:b1:comparative-planetology:moi}{moment-of-inertia factor}
(question 15) and the surface gravity (question 17, against the measured mean
\qty{9.82}{m.s^{-2}}) are checked against measurements of the Earth; the
gravity at the core and the pressures can only be compared with PREM, itself
a model.

\textbf{24.} The core holds \textbf{\qty[text-series-to-math=true]{32.5}{\percent}}
of the Earth's mass.
\end{solution}
